\documentclass[11pt]{article}
\usepackage[a4paper,margin=25mm]{geometry}
\usepackage{amsmath,amssymb,amsthm,mathtools}
\usepackage[colorlinks=true,linkcolor=blue,urlcolor=blue]{hyperref}
\newtheorem{theorem}{Theorem}
\newtheorem{lemma}{Lemma}
\newcommand{\Tr}{\operatorname{Tr}}
\newcommand{\A}{\mathcal A}
\newcommand{\R}{\mathcal R}
\title{Order monotonicity of anticoherence measures}
\author{}\date{24 November 2026}
\begin{document}
\maketitle
\noindent
The fidelity-based total and quantum anticoherence measures are nonincreasing in the order. We give a general proof and a Lean certificate. A single $W$ state provides explicit counterexamples for the purity, von Neumann and Hilbert--Schmidt constructions.\footnote{The fidelity proof is checked in Lean~4.24.0, with no admitted proofs or additional mathematical axioms. Exact versions, formal statements and reproduction instructions are in \texttt{lean/README.md}.}

\begin{theorem}[Order monotonicity]
For every symmetric $N$-qubit state $\rho$ and $2\le t<N$,
\begin{equation}\label{eq:main}
 \A_t^{T,F}(\rho)\le \A_{t-1}^{T,F}(\rho),\qquad
 \A_t^{Q,F}(\rho)\le \A_{t-1}^{Q,F}(\rho).
\end{equation}
Here
\[
 \A_t^{T,F}(\rho)=\frac{(\Tr\sqrt{\rho_t})^2-1}{t},\qquad
 \A_t^{Q,F}(\rho)=\inf_{\rho=\sum_a p_a|\psi_a\rangle\langle\psi_a|}
 \sum_a p_a\A_t^{T,F}(\psi_a).
\]
The total inequality also holds at $t=N$ if the same total formula is used there.
\end{theorem}

\noindent The main ingredient is a statement for dimension-changing positive maps. For a density matrix $X$ of dimension $d\ge2$, define
\[
 \R_d(X)=\frac{(\Tr\sqrt X)^2-1}{d-1}.
\]

\begin{lemma}[Dimension reduction]\label{lem:map}
Let $\Phi:M_{n+1}\to M_n$, $n\ge2$, be positive and trace preserving, and suppose
\begin{equation}\label{eq:scalar}
 \Phi(I_{n+1})=\frac{n+1}{n}I_n.
\end{equation}
Then, for every density matrix $X$,
\begin{equation}\label{eq:general}
 \R_n(\Phi(X))\ge \R_{n+1}(X).
\end{equation}
Complete positivity is not needed.
\end{lemma}

\begin{proof}
Write the eigenvalues of $X$ as
$\lambda_1\ge\cdots\ge\lambda_{n+1}\ge0$, and set
\begin{equation}\label{eq:mu}
 \mu_i=\frac{n-i+1}{n}\lambda_i+\frac{i}{n}\lambda_{i+1},\qquad 1\le i\le n.
\end{equation}
We first show that $\lambda(\Phi(X))\prec\mu$. The entries of $\mu$ sum to one and are decreasing, since
\[
 n(\mu_i-\mu_{i+1})=(n-i+1)(\lambda_i-\lambda_{i+1})
 +(i+1)(\lambda_{i+1}-\lambda_{i+2})\ge0.
\]
For every rank-$k$ projection $P$, $1\le k<n$, the adjoint satisfies
\[
 0\le\Phi^*(P)\le I_{n+1},\qquad
 \Tr\Phi^*(P)=\frac{k(n+1)}n=k+\frac kn.
\]
Consequently,
\begin{equation}\label{eq:kyfan}
 \Tr[P\Phi(X)]\le\sum_{i=1}^k\lambda_i+\frac kn\lambda_{k+1}
 =\sum_{i=1}^k\mu_i.
\end{equation}
For completeness, the inequality follows by diagonalizing $X$: the diagonal entries of $\Phi^*(P)$ lie in $[0,1]$ and sum to $k+k/n$. Their scalar product with decreasing $\lambda_i$ is maximized by placing weights $1$ on the first $k$ entries and $k/n$ on the next one. Maximizing the left side of \eqref{eq:kyfan} over $P$ proves the asserted majorization. Concavity of the square root therefore gives
\begin{equation}\label{eq:spectral}
 (\Tr\sqrt{\Phi(X)})^2\ge\left(\sum_{i=1}^n\sqrt{\mu_i}\right)^2.
\end{equation}

We next prove the scalar estimate
\begin{equation}\label{eq:scalarineq}
 \left(\sum_{i=1}^n\sqrt{\mu_i}\right)^2-1
 \ge\frac{n-1}{n}
 \left[\left(\sum_{i=1}^{n+1}\sqrt{\lambda_i}\right)^2-1\right].
\end{equation}
Put $x_i=\sqrt{\lambda_i}$, $a_i=(n-i+1)/n$ and $b_i=i/n$. For $i<j$, the weighted Cauchy--Schwarz inequality yields
\begin{equation}\label{eq:pair}
 \sqrt{\mu_i\mu_j}\ge
 a_jx_ix_j+(b_j-b_i)x_ix_{j+1}+b_ix_{i+1}x_{j+1}.
\end{equation}
Indeed, apply it to the vectors
\[
 (\sqrt{a_j}x_i,\sqrt{b_j-b_i}x_i,\sqrt{b_i}x_{i+1}),\quad
 (\sqrt{a_j}x_j,\sqrt{b_j-b_i}x_{j+1},\sqrt{b_i}x_{j+1}),
\]
whose squared norms are $\mu_i$ and $\mu_j$.
Summing \eqref{eq:pair} over $1\le i<j\le n$ gives
\begin{equation}\label{eq:sum}
 \sum_{i<j\le n}\sqrt{\mu_i\mu_j}
 \ge\frac{n-1}{n}\sum_{k<\ell\le n+1}x_kx_\ell.
\end{equation}
To check the coefficient of $x_kx_\ell$, the three terms in \eqref{eq:pair} contribute respectively
$(n-\ell+1)/n$, $(\ell-k-1)/n$, and $(k-1)/n$.
An absent term has zero coefficient, including adjacent indices and boundary indices. Their sum is $(n-1)/n$.
Since $\sum_i\mu_i=\sum_i x_i^2=1$, \eqref{eq:sum} is exactly \eqref{eq:scalarineq}. Combining it with \eqref{eq:spectral} proves \eqref{eq:general}.
\end{proof}

\begin{proof}[Proof of the theorem]
The partial trace of one qubit, restricted to the symmetric sector, is a completely positive trace-preserving map
\[
 T_t:\mathcal B(\operatorname{Sym}^t\mathbb C^2)
 \longrightarrow\mathcal B(\operatorname{Sym}^{t-1}\mathbb C^2).
\]
Rotational covariance and irreducibility imply
\[
 T_t(I_{t+1})=\frac{t+1}{t}I_t.
\]
The scalar factor also follows by taking the trace. Applying Lemma~\ref{lem:map} with $n=t$, $X=\rho_t$, and $T_t(\rho_t)=\rho_{t-1}$ proves the total inequality.

For every pure-state decomposition of $\rho$,
\[
 \sum_a p_a\A_t^{T,F}(\psi_a)
 \le\sum_a p_a\A_{t-1}^{T,F}(\psi_a).
\]
Taking the infimum over the same set of decompositions on both sides proves the quantum inequality. Equivalently, use a minimizing sequence for the order-$(t-1)$ roof.
\end{proof}

\paragraph{Equivalent fidelity contraction.}
With squared Uhlmann fidelity $F$ and $\tau=T_t(\sigma)$, the general lemma gives the sharp bound
\[
 1-F\left(\tau,\frac{I_t}{t}\right)
 \le\left(1-\frac1{t^2}\right)
 \left[1-F\left(\sigma,\frac{I_{t+1}}{t+1}\right)\right].
\]
Pure coherent states saturate it, so its coefficient cannot be decreased.


\clearpage
\section*{Simple counterexamples for the other three measures}
A single Dicke state suffices. Let
\[
 |W_N\rangle=\frac1{\sqrt N}\sum_{i=1}^N
 |0\cdots 0\,1_i\,0\cdots0\rangle.
\]
Its $t$-qubit marginal is
\[
 \rho_t=\left(1-\frac tN\right)|0\rangle^{\otimes t}\langle0|^{\otimes t}
       +\frac tN|W_t\rangle\langle W_t|,
\]
with nonzero eigenvalues $1-t/N$ and $t/N$. On the pure full state, each quantum convex roof equals its total pure-state functional. Hence
\begin{align*}
 \A_t^{Q,\mathcal P}(W_N)&=\frac{2(t+1)(N-t)}{N^2},\\
 \A_t^{Q,\mathrm{HS}}(W_N)&=1-\sqrt{1-\frac{2(t+1)(N-t)}{N^2}},\\
 \A_t^{Q,\mathrm{vN}}(W_N)&=\frac{h(t/N)}{\log(t+1)},
 \qquad h(x)=-x\log x-(1-x)\log(1-x).
\end{align*}
Thus the same pure state $|W_{14}\rangle$ violates order monotonicity for all three:
\begin{center}
\renewcommand{\arraystretch}{1.35}
\begin{tabular}{lcc}
Construction & $\A_1^Q$ & $\A_2^Q$\\\hline
Purity & $13/49$ & $18/49$\\
Hilbert--Schmidt & $1/7$ & $1-\sqrt{31}/7$\\
von Neumann & $h(1/14)/\log2$ & $h(1/7)/\log3$\\
 & $0.371232\ldots$ & $0.373304\ldots$
\end{tabular}
\end{center}
The second-order value is strictly larger in every row. For the von Neumann comparison, exact rational interval bounds give
\[
 0.37123232<\A_1^{Q,\mathrm{vN}}(W_{14})<0.37123234,
 \qquad
 0.37330395<\A_2^{Q,\mathrm{vN}}(W_{14})<0.37330398.
\]
These are certified using rational arithmetic and the expansion, for $z=(x-1)/(x+1)\in(0,1)$,
\[
 \log x=2\sum_{k=0}^{M-1}\frac{z^{2k+1}}{2k+1}+R_M,
 \qquad 0<R_M\le\frac{2z^{2M+1}}{(2M+1)(1-z^2)}.
\]
The included script \texttt{verify\_counterexamples.py} uses $M=64$; its comparisons do not use floating-point arithmetic.

For purity and HS, an even smaller example is $|W_5\rangle$:
\[
 \A_1^{Q,\mathcal P}=\frac{16}{25}<\frac{18}{25}=\A_2^{Q,\mathcal P},
 \qquad
 \A_1^{Q,\mathrm{HS}}=\frac25<1-\frac{\sqrt7}{5}=\A_2^{Q,\mathrm{HS}}.
\]
In contrast, fidelity satisfies
$\A_t^{Q,F}(W_N)=2\sqrt{(N-t)/t}/N$, which decreases with $t$, consistently with the general theorem.

\end{document}
